%% Conclusion: What the Program Has Shown

This thesis began with a question that felt philosophical: what must
an observer be, for coherent probabilistic experience to be possible?
It ends with an algorithm.  The distance between those two points is
the measure of what the program has established.

\section*{The arc}

Chapter~\ref{chap:paper-minus-one} showed that an observer whose
discriminative acts satisfy finite Boolean closure, and who is
committed to acting coherently toward a world that exceeds their
current distinctions, is forced to possess a $\sigma$-algebra.  Not
as an assumption.  As the form that coherent commitment must take.

Chapter~\ref{chap:paper-zero} showed that when a compatible family
of such observers --- queries organized by a refinement order ---
satisfies collective exhaustion and a realizability condition on the
projective limit, their local commitments assemble into a unique
global probability measure $P$ on the realization space $\Omega$.
Nothing is hidden: $P$ is entirely determined by the observable
interface.

Chapter~\ref{chap:paper-one} showed that the predictive content of
each query factors through a minimal sufficient statistic: the
minimal predictive query $Q_*$.  Prediction is structured, not
arbitrary.

Chapter~\ref{chap:paper-two} showed that as time advances, the
predictive structure evolves as a semigroup.  The dynamics are
coherent across time in the same way that the queries are coherent
across refinement levels.

Chapter~\ref{chap:paper-three} showed that this semigroup can be
instantiated via delay queries and approximated from finite data.
The abstract program becomes computable.

\section*{What was derived, what was assumed}

The program is honest about the distinction.  Three things were
derived: the $\sigma$-algebra (from coherent commitment and the
extension criterion, given collective exhaustion), $\sigma$-additivity
(from tightness or perfectness, for the architectural cases), and the
predictive semigroup (from observable compatibility alone).

Two things remain as structural hypotheses: upper-directedness
(SP4: assumed, with an operational motivation but no derivation),
and the general case of SP2 beyond compact and standard Borel
outcome spaces.

Realizability of $\Omega$ (SP3) is partially resolved.
Independence is established: a system can satisfy all three other
hypotheses while failing realizability, the obstruction being
a non-surjective refinement map.
Surjectivity of all refinement maps is therefore a \emph{necessary}
condition.  It is not abstractly sufficient: the general inverse
limit theorem requires compactness of outcome spaces.
For the compact Hausdorff setting, it is now a theorem: continuous
surjective refinement maps between compact T2 outcome spaces force
realizability (formalized as
\texttt{TopologicalQuerySystem.evalSurjective\_of\_compact\_t2},
zero sorrys, via Tychonoff compactness and a partial-coherence
intersection argument).  This covers the programme's concrete
systems --- delay query systems and finite-outcome systems.
The abstract sufficient condition in full generality (without
compactness) remains open.

These are honest gaps.  The program does not claim to derive
everything from nothing.  It claims to derive more than was
previously known to be derivable, and to identify precisely what
the remaining assumptions are and what they mean.

\section*{The open frontier}

The CE irreducibility result of
Chapter~\ref{chap:paper-minus-one} --- that no structural condition
on a query system can force collective exhaustion --- is a proved
boundary, not an open question.  The finite-cofinite counterexample
on $\mathbb{Q}$ witnesses the particular case; the Łoś ultraproduct
argument establishes that no first-order condition on the query
system suffices.  The gap between coherence and probability is exact.

What this means for the chain:
\[
\text{bounded discriminability}
\;\to\;
\text{refinement}
\;\to\;
\text{coherence}
\;\xrightarrow{\;\text{CE}\;}\;
\text{probability}
\;\to\;
\text{prediction}
\;\to\;
\text{dynamics}
\;\to\;
\text{computation.}
\]
The arrow from coherence to probability requires collective
exhaustion --- an irreducible commitment that cannot be derived
from structural conditions alone.  Every other arrow is forced.
CE names precisely what must be assumed and shows it cannot be
smaller.  That is not a failure.  It is the program's deepest
result: the exact location and minimum size of the commitment
probability requires.

The remaining open question is the philosophical status of CE
as a primitive.  CE is a commitment about the completed infinite
made by a finite observer: the willingness to be corrected by
what the infinite reveals, even knowing one can never reach it.
Whether that commitment can be \emph{justified} from within any
formal system is suspected to be undecidable in character ---
the query-system analogue of the axiom of infinity.  That
question is earmarked, not abandoned.
